% SCP10, paper_v3.tex:2832--2940, especially 2861--2877 and 2896--2923.
% The available tc-lattice and tc-blocked-lattice artwork was inspected.
% The unavailable renorm-as-diff-doubles artwork is not reconstructed here.
% Declaration links are added only after the corresponding statements pass.
\paragraph{The clockwise color-difference tensor.}
Let $G$ be a finite group. At one blocked site, place virtual colors
$(p,q,r,s)$ on the north, east, south and west legs, respectively.
The physical tuple $(a,b,c,d)$ is ordered northeast, southeast, southwest,
northwest. The ordered differences follow
\cite[Equation~(7.10) and its clockwise-convention footnote]{Schuch2010PEPS}.
The earlier checkerboard drawings specify the incidence and blocking
pattern. The physical orientation and multiplication order here are fixed
by the clockwise convention of the displayed tensor; no arrow-by-arrow
identification with those earlier drawings is imposed.

\begin{definition}[The clockwise blocked tensor]
    \label{def:peps_quantum_double_k_tensor}
    \lean{TNLean.PEPS.quantumDoubleKSpins,TNLean.PEPS.quantumDoubleKTensor,
        TNLean.PEPS.siteMap_quantumDoubleKTensor_apply}
    \leanok
    Define the color-difference map $F:G^4\to G^4$ and the tensor $K$ by
    \begin{align}
        F(p,q,r,s)         & =(pq^{-1},qr^{-1},rs^{-1},sp^{-1}), \\
        K(p,q,r,s;a,b,c,d) & =\delta_{(a,b,c,d),F(p,q,r,s)}.
        \label{eq:peps_qd_k_tensor}
    \end{align}
    Its linear map is $\mathsf K\ket{p,q,r,s}=\ket{F(p,q,r,s)}$.
    No normalization is included in $K$.
\end{definition}
% Formula: K(p,q,r,s;a,b,c,d)=delta_((a,b,c,d),F(p,q,r,s)).
% Ink: virtual N,E,S,W = p,q,r,s; physical NE,SE,SW,NW = a,b,c,d.
% All eight legs are open; no repeated label is summed in this picture.
% External boundary signature: (virtual,physical)=(4,4), alphabet G per leg.
\begin{tenkzequation}
    \tenkzkernel
    \begin{tenkz}[rows={wire},cols=1,bonds=none]
        \tn[at=(1,1),name=qdK,skin=box,
            ports={90:virtual:$p$,0:virtual:$q$,270:virtual:$r$,180:virtual:$s$,
                    45:physical:$a$,315:physical:$b$,225:physical:$c$,135:physical:$d$}]{K}
    \end{tenkz}
\end{tenkzequation}

\begin{theorem}[Comparison of physical orientation conventions]
    \label{thm:peps_quantum_double_k_orientation}
    \lean{TNLean.PEPS.quantumDoubleDualToK,
        TNLean.PEPS.quantumDoubleDualToK_spins,
        TNLean.PEPS.quantumDoubleKTensor_eq_dual}
    \leanok
    \uses{def:peps_quantum_double_k_tensor}
    For the dual convention
    $F_{\rm dual}(p,q,r,s)=(ps^{-1},qp^{-1},qr^{-1},rs^{-1})$, the
    physical basis permutation
    $J(y_1,y_2,y_3,y_4)=(y_2^{-1},y_3,y_4,y_1^{-1})$ satisfies
    $J\circ F_{\rm dual}=F$.
\end{theorem}
\begin{proof}\leanok
    The first and fourth components follow from
    $(qp^{-1})^{-1}=pq^{-1}$ and $(ps^{-1})^{-1}=sp^{-1}$.
    The other two components retain their original order.
\end{proof}
% J is a local physical coordinate permutation of two blocked conventions;
% it is not a unitary renormalization of the unblocked tensor T.

\begin{definition}[The local flatness projector]
    \label{def:peps_quantum_double_local_projector}
    \lean{TNLean.PEPS.quantumDoubleKHolonomy,
        TNLean.PEPS.quantumDoubleKLocalProjector,TNLean.PEPS.quantumDoubleKLocalTerm,
        TNLean.PEPS.quantumDoubleKLocalProjector_mulVec}
    \leanok
    \uses{def:peps_quantum_double_k_tensor}
    On the full four-spin space $\C[G^4]$, set
    \begin{align}
        P_0 & =\sum_{abcd=1}\ketbra{a,b,c,d}{a,b,c,d},
            & h_0                                      & =I-P_0.
        \label{eq:peps_qd_local_projector}
    \end{align}
    The product is ordered as displayed; $G$ need not be abelian.
\end{definition}

\begin{theorem}[The actual local tensor image]
    \label{thm:peps_quantum_double_local_image}
    \lean{TNLean.PEPS.quantumDoubleKSpins_range,
        TNLean.PEPS.quantumDoubleKHolonomy_spins,
        TNLean.PEPS.quantumDoubleKColors,TNLean.PEPS.quantumDoubleKSpins_colors,
        TNLean.PEPS.siteMap_quantumDoubleKTensor_single,
        TNLean.PEPS.quantumDoubleKLocalProjector_isStarProjection,
        TNLean.PEPS.quantumDoubleKLocalTerm_isStarProjection,
        TNLean.PEPS.quantumDoubleKLocalProjector_siteMap,
        TNLean.PEPS.range_siteMap_quantumDoubleKTensor,
        TNLean.PEPS.range_siteMap_quantumDoubleKTensor_eq_ker_localTerm}
    \leanok
    \uses{def:peps_quantum_double_local_projector}
    The operators $P_0$ and $h_0$ are orthogonal projectors, with
    $P_0\mathsf K=\mathsf K$ and $\operatorname{ran}\mathsf K=\ker h_0$.
    \begin{align}
        \operatorname{ran}\mathsf K
         & =\operatorname{ran}P_0
        =\operatorname{span}\{\ket{a,b,c,d}:abcd=1\}.
        \label{eq:peps_qd_local_image}
    \end{align}
\end{theorem}
\begin{proof}\leanok
    The ordered product of the four components of $F(p,q,r,s)$ is one.
    Conversely, if $abcd=1$, the coloring
    $(1,a^{-1},(ab)^{-1},(abc)^{-1})$ maps to $(a,b,c,d)$.
    Applying $\mathsf K$ to that virtual basis vector therefore gives the
    indicated physical basis vector. This proves the two range inclusions.
    The diagonal entries of $P_0$ are zero or one, so $P_0^2=P_0=P_0^\dagger$;
    its complement has the same properties.
    Idempotence gives $\operatorname{ran}P_0=\ker(I-P_0)$.
\end{proof}

\paragraph{One shared bond.}
The east color of a left tensor and the west color of a right tensor are
one summed index $x$. With six exterior colors fixed, the actual open
contraction is
\begin{align}
    \Psi_{p,r,s;t,v,w}
     & =\sum_{x\in G}\ket{F(p,x,r,s)}\otimes\ket{F(t,v,w,x)}.
    \label{eq:peps_qd_bond_contraction}
\end{align}
% Formula: Psi=sum_x |F(p,x,r,s)> tensor |F(t,v,w,x)>.
% Ink: the sole internal virtual wire is q(left)=s(right)=x, summed once.
% The six other virtual colors p,r,s,t,v,w are open boundary arguments.
% Each physical leg is a complete G^4 tuple, not one selected corner spin.
% External boundary signature: (virtual,physical)=(6,2).
\begin{tenkzequation}
    \tenkzkernel
    \begin{tenkz}[rows={wire},cols=2,bonds=none]
        \tn[at=(1,1),name=qdLeft,skin=box,
            ports={90:virtual:$p$,0:virtual,270:virtual:$r$,180:virtual:$s$,
                    135:physical:$\sigma_L$}]{K_L}
        \tn[at=3 e of qdLeft,name=qdRight,skin=box,
            ports={90:virtual:$t$,0:virtual:$v$,270:virtual:$w$,180:virtual,
                    45:physical:$\sigma_R$}]{K_R}
        \tnwire[name=qdShared]{qdLeft.0}{qdRight.180}
        \tnmark[form=label,label pos=n]{on qdShared 0.5}{$x$}
    \end{tenkz}
\end{tenkzequation}

\begin{definition}[Coherent action at an east--west bond]
    \label{def:peps_quantum_double_bond_average}
    \lean{TNLean.PEPS.quantumDoubleKBondContraction,
        TNLean.PEPS.quantumDoubleKBondContraction_eq,
        TNLean.PEPS.quantumDoubleKBondLeft,TNLean.PEPS.quantumDoubleKBondRight,
        TNLean.PEPS.quantumDoubleKBondAction,TNLean.PEPS.quantumDoubleKBondPermutation,
        TNLean.PEPS.quantumDoubleKBondMatrix,TNLean.PEPS.quantumDoubleKBondMatrix_mulVec,
        TNLean.PEPS.quantumDoubleKBondAverage,
        TNLean.PEPS.quantumDoubleKBondTerm}
    \leanok
    \uses{def:peps_quantum_double_k_tensor}
    For $u\in G$, define the permutation operator $V_u$ on
    $\C[G^4]\otimes\C[G^4]$ by
    \begin{align}
        V_u\bigl(\ket{a,b,c,d}\otimes\ket{e,f,g,h}\bigr)
            & =\ket{au^{-1},ub,c,d}\otimes\ket{e,f,gu^{-1},uh},           \\
        A_e & =\frac1{|G|}\sum_{u\in G}V_u,
            & h_e                                               & =I-A_e.
        \label{eq:peps_qd_bond_average}
    \end{align}
    The four changed spins carry $R_u,L_u,R_u,L_u$, respectively, where
    $L_u\ket z=\ket{uz}$ and $R_u\ket z=\ket{zu^{-1}}$.
    These are the incoming/outgoing spin actions of
    \cite[Section~7, ``The double models'']{Schuch2010PEPS}.
\end{definition}

\begin{theorem}[The normalized bond projector fixes the contraction]
    \label{thm:peps_quantum_double_bond_fixed}
    \lean{TNLean.PEPS.quantumDoubleKBondAction_mul,
        TNLean.PEPS.quantumDoubleKBondAction_one,
        TNLean.PEPS.quantumDoubleKSpins_bondLeft,TNLean.PEPS.quantumDoubleKSpins_bondRight,
        TNLean.PEPS.quantumDoubleKBondMatrix_single,
        TNLean.PEPS.quantumDoubleKBondAverage_isStarProjection,
        TNLean.PEPS.quantumDoubleKBondTerm_isStarProjection,
        TNLean.PEPS.quantumDoubleKBondAverage_mulVec_eq_self_iff,
        TNLean.PEPS.quantumDoubleKBondContraction_invariant,
        TNLean.PEPS.quantumDoubleKBondAverage_contraction}
    \leanok
    \uses{def:peps_quantum_double_bond_average}
    The operators satisfy $V_uV_v=V_{uv}$ and $V_u^\dagger=V_{u^{-1}}$.
    The average $A_e$ and its complement $h_e$ are orthogonal projectors.
    For every physical vector $v$, the equality $A_ev=v$ holds exactly
    when $V_uv=v$ for all $u\in G$.
    For every choice of the six exterior colors,
    \begin{align}
        V_u\Psi_{p,r,s;t,v,w} & =\Psi_{p,r,s;t,v,w},        \\
        A_e\Psi_{p,r,s;t,v,w} & =\Psi_{p,r,s;t,v,w},
                              & h_e\Psi_{p,r,s;t,v,w} & =0.
        \label{eq:peps_qd_bond_fixed}
    \end{align}
\end{theorem}
\begin{proof}\leanok
    Composition follows from $(uv)^{-1}=v^{-1}u^{-1}$; the inverse
    permutation is $V_{u^{-1}}$. In the double sum for $A_e^2$, each
    product $uv$ occurs exactly $|G|$ times. Inversion permutes the sum
    for $A_e^\dagger$, proving both projector identities.
    The image of the group average is its invariant subspace.
    On each summand of (\ref{eq:peps_qd_bond_contraction}), $V_u$ replaces
    the common virtual color $x$ by $ux$ in both tensors. The bijection
    $x\mapsto ux$ preserves the sum. Averaging the resulting identities
    gives (\ref{eq:peps_qd_bond_fixed}).
\end{proof}

The factor $|G|^{-1}$ in (\ref{eq:peps_qd_bond_average}) is necessary.
For $S_e=\sum_uV_u$, one has $S_e^2=|G|S_e$; on an invariant vector,
$(I-S_e)\psi=(1-|G|)\psi$. Thus the unnormalized sum printed in
\cite[Section~7, ``The double models'']{Schuch2010PEPS} must be divided
by $|G|$ to obtain the projector complement that annihilates the state.
% Local correction and remaining scope are recorded in
% docs/paper-gaps/scp10_quantum_double_local_hamiltonian.tex.

\begin{theorem}[The actual two-block tensor image]
    \label{thm:peps_quantum_double_bond_image}
    \lean{TNLean.PEPS.quantumDoubleKPairLocalProjector,
        TNLean.PEPS.quantumDoubleKHolonomy_bondLeft,
        TNLean.PEPS.quantumDoubleKHolonomy_bondRight,
        TNLean.PEPS.quantumDoubleKPairLocalProjector_isStarProjection,
        TNLean.PEPS.quantumDoubleKPairLocalProjector_commute_bondAverage,
        TNLean.PEPS.quantumDoubleKBondSupportProjector,
        TNLean.PEPS.quantumDoubleKBondSupportProjector_isStarProjection,
        TNLean.PEPS.quantumDoubleKPairLocalProjector_contraction,
        TNLean.PEPS.quantumDoubleKBondSupportProjector_contraction,
        TNLean.PEPS.quantumDoubleKBondBoundaryColors,
        TNLean.PEPS.quantumDoubleKBondBoundaryColors_spins,
        TNLean.PEPS.quantumDoubleKBondAverage_single_eq_contraction,
        TNLean.PEPS.range_quantumDoubleKBondContraction,
        TNLean.PEPS.mem_range_quantumDoubleKBondContraction_iff}
    \leanok
    \uses{def:peps_quantum_double_local_projector,
        thm:peps_quantum_double_bond_fixed}
    Let $\mathsf C\ket\beta=\Psi_\beta$ be the open contraction map in
    (\ref{eq:peps_qd_bond_contraction}), and let $P_{00}=P_0\otimes P_0$.
    Then $P_{00}A_e=A_eP_{00}$, the product $Q_e=A_eP_{00}$ is an
    orthogonal projector, and
    \begin{align}
        \operatorname{ran}\mathsf C & =\operatorname{ran}Q_e.
        \label{eq:peps_qd_bond_image}
    \end{align}
    Equivalently, a physical vector $v$ belongs to this actual contraction
    image exactly when $P_{00}v=v$ and $V_uv=v$ for every $u\in G$.
\end{theorem}
\begin{proof}\leanok
    Each bond action preserves both ordered local products, so its
    permutation matrix commutes with $P_{00}$. Hence the average commutes
    with $P_{00}$ and their product is an orthogonal projector.
    Each contraction column is fixed by both factors, proving one inclusion.
    For the reverse inclusion, take
    $\xi=(a,b,c,d;e,f,g,h)$ with $abcd=efgh=1$ and set
    $\beta(\xi)=(a,b^{-1},(bc)^{-1};efg,fg,g)$. At shared color one,
    these boundary colors produce precisely $\ket\xi$. Varying the shared
    color gives
    \begin{align}
        A_e\ket\xi & =|G|^{-1}\Psi_{\beta(\xi)}.
    \end{align}
    The projector $P_{00}$ kills every other computational basis vector.
    Expand an arbitrary vector in that basis to obtain the second inclusion.
\end{proof}

\paragraph{The hole between four blocked tensors.}
Place four copies of $K$ at northwest, northeast, southeast and southwest.
Let their four shared colors, starting at the northern bond, be $x,y,z,w$
on the north, west, south and east bonds, respectively. Fix the eight
exterior colors $\beta=(p,s,t,v,k,l,m,n)$ and define
\begin{align}
    \Phi_\beta
     & =\sum_{x,y,z,w\in G}
    \ket{F(p,x,y,s)}\otimes\ket{F(t,v,w,x)}\nonumber                 \\
     & \hspace{2.5em}\otimes\ket{F(w,k,l,z)}\otimes\ket{F(y,z,m,n)}.
    \label{eq:peps_qd_hole_contraction}
\end{align}
The tensor factors are ordered northwest, northeast, southeast, southwest.
% Formula: Phi_beta is the literal four-K sum in eq:peps_qd_hole_contraction.
% Ink: NW.E=NE.W=x; NW.S=SW.N=y; SW.E=SE.W=z; NE.S=SE.N=w.
% The four internal virtual colors are each summed once; exterior colors
% p,s,t,v,k,l,m,n are independent, unsummed boundary arguments.
% Physical sigma_NW,sigma_NE,sigma_SE,sigma_SW are complete G^4 tuples.
% The selected hole spins are NW.b,SW.a,SE.d,NE.c, in that product order.
% External boundary signature: (virtual,physical)=(8,4).
\begin{tenkzequation}
    \tenkzkernel
    \begin{tenkz}[rows={wire,wire},cols=2,bonds=none]
        \tn[at=(1,1),name=qdNW,skin=box,
            ports={90:virtual:$p$,0:virtual,270:virtual,180:virtual:$s$,
                    135:physical:$\sigma_{NW}$}]{K}
        \tn[at=3 e of qdNW,name=qdNE,skin=box,
            ports={90:virtual:$t$,0:virtual:$v$,270:virtual,180:virtual,
                    45:physical:$\sigma_{NE}$}]{K}
        \tn[at=3 s of qdNE,name=qdSE,skin=box,
            ports={90:virtual,0:virtual:$k$,270:virtual:$l$,180:virtual,
                    315:physical:$\sigma_{SE}$}]{K}
        \tn[at=3 s of qdNW,name=qdSW,skin=box,
            ports={90:virtual,0:virtual,270:virtual:$m$,180:virtual:$n$,
                    225:physical:$\sigma_{SW}$}]{K}
        \tnwire[name=qdNorth]{qdNW.0}{qdNE.180}
        \tnwire[name=qdWest]{qdNW.270}{qdSW.90}
        \tnwire[name=qdSouth]{qdSW.0}{qdSE.180}
        \tnwire[name=qdEast]{qdNE.270}{qdSE.90}
        \tnmark[form=label,label pos=n]{on qdNorth 0.5}{$x$}
        \tnmark[form=label,label pos=w]{on qdWest 0.5}{$y$}
        \tnmark[form=label,label pos=s]{on qdSouth 0.5}{$z$}
        \tnmark[form=label,label pos=e]{on qdEast 0.5}{$w$}
    \end{tenkz}
\end{tenkzequation}

\begin{definition}[The four-block plaquette projector]
    \label{def:peps_quantum_double_hole_projector}
    \lean{TNLean.PEPS.QuantumDoubleKPlaquetteConfig,
        TNLean.PEPS.quantumDoubleKPlaquetteSpins,
        TNLean.PEPS.quantumDoubleKPlaquetteContraction,
        TNLean.PEPS.quantumDoubleKPlaquetteHolonomy,
        TNLean.PEPS.quantumDoubleKPlaquetteProjector,
        TNLean.PEPS.quantumDoubleKPlaquetteTerm}
    \leanok
    \uses{def:peps_quantum_double_k_tensor}
    For $\sigma=(\sigma_{NW},\sigma_{NE},\sigma_{SE},\sigma_{SW})$, set
    \begin{align}
        H_\square(\sigma)   & =\sigma_{NW,b}\sigma_{SW,a}
        \sigma_{SE,d}\sigma_{NE,c},                                                    \\
        P_\square\ket\sigma & =\delta_{H_\square(\sigma),1}\ket\sigma,
                            & h_\square                                & =I-P_\square.
        \label{eq:peps_qd_hole_projector}
    \end{align}
    Both operators act on the full sixteen-spin space.
\end{definition}

\begin{theorem}[Plaquette flatness of the actual four-block contraction]
    \label{thm:peps_quantum_double_hole_fixed}
    \lean{TNLean.PEPS.quantumDoubleKPlaquetteHolonomy_spins,
        TNLean.PEPS.quantumDoubleKPlaquetteProjector_isStarProjection,
        TNLean.PEPS.quantumDoubleKPlaquetteTerm_isStarProjection,
        TNLean.PEPS.quantumDoubleKPlaquetteContraction_eq,
        TNLean.PEPS.quantumDoubleKPlaquetteProjector_contraction}
    \leanok
    \uses{def:peps_quantum_double_hole_projector}
    The operators $P_\square$ and $h_\square$ are orthogonal projectors.
    For every exterior coloring $\beta$, one has
    $P_\square\Phi_\beta=\Phi_\beta$ and $h_\square\Phi_\beta=0$.
\end{theorem}
\begin{proof}\leanok
    Each summand in (\ref{eq:peps_qd_hole_contraction}) has selected spins
    $(xy^{-1},yz^{-1},zw^{-1},wx^{-1})$. Their product is one by
    adjacent cancellation, without commuting any two factors.
    The diagonal projector therefore fixes each summand and their sum.
    Its entries are zero or one, proving the orthogonal-projector assertions.
\end{proof}

\begin{theorem}[Commutation at the northern bond]
    \label{thm:peps_quantum_double_hole_commute}
    \lean{TNLean.PEPS.quantumDoubleKPlaquetteTopBond,
        TNLean.PEPS.quantumDoubleKPlaquetteHolonomy_topBond,
        TNLean.PEPS.quantumDoubleKPlaquetteHolonomy_topBond_eq_one_iff,
        TNLean.PEPS.quantumDoubleKPlaquetteProjector_commute_topBond,
        TNLean.PEPS.quantumDoubleKPlaquetteProjector_commute_topBondAverage}
    \leanok
    \uses{def:peps_quantum_double_bond_average,
        def:peps_quantum_double_hole_projector}
    Let $W_u$ act by $V_u$ on the northwest and northeast physical tuples
    and by the identity on the southeast and southwest tuples.
    The hole projector commutes with every $W_u$ and with
    $|G|^{-1}\sum_uW_u$.
\end{theorem}
\begin{proof}\leanok
    The northern bond changes the ordered holonomy by conjugation:
    $H_\square(W_u\sigma)=uH_\square(\sigma)u^{-1}$.
    Consequently the product-one indicator is unchanged on the entire
    physical basis, including configurations with nontrivial holonomy.
    The diagonal projector therefore commutes with the permutation and
    with its normalized average.
\end{proof}

% Scope: exact local and east--west two-block terms, and a four-block hole.
% These statements do not identify a full-lattice common kernel, prove
% commutation of all placed terms, or construct the unblocked T-to-K unitary.
